\chapter{度量空间的拓扑}\label{ch:b2:metric}

实直线（本科一年级分册）上的拓扑，几乎不改一字地
推广到任意赋有距离的集合上。
其收益巨大：函数列、矩阵、曲线——
都成为\omterm{def:b2:metric:def}{度量空间}中的点，本章证明的三大支柱
——\omterm{def:b2:metric:complete}{完备}性（连同巴拿赫不动点定理）、紧性、
连通性——对它们一致地适用。本章是
全书分析部分的骨架。

\section{度量空间}

\begin{definition}\label{def:b2:metric:def}
称\emph{度量空间}\index{度量空间}为集合 $X$，并在其上给定映射 $d
\colon X \times X \to \R_+$，使对所有 $x, y, z$：
\[
d(x,y) = 0 \iff x = y,
\qquad
d(x,y) = d(y,x),
\qquad
d(x,z) \leq d(x,y) + d(y,z).
\]
球：开球 $B(a, r) = \{x : d(a,x) < r\}$，闭球 $\overline B(a,r) =
\{x : d(a,x) \leq r\}$。子集 $A \subseteq X$ 在诱导距离下
构成度量空间。
\end{definition}

\begin{example}\label{ex:b2:metric:examples}
$\R$ 上取 $\abs{x - y}$；$\R^n$ 上可取下列任一距离：
\[
d_1(x,y) = \sum_i \abs{x_i - y_i},
\quad
d_2(x,y) = \Bigl(\sum_i (x_i - y_i)^2\Bigr)^{1/2},
\quad
d_\infty(x,y) = \max_i \abs{x_i - y_i};
\]
\omterm{def:b2:metric:continuity}{连续}函数构成的集合 $C(\intcc{a}{b})$，赋以
\emph{上确界距离} $d_\infty(f, g) = \sup_{\intcc{a}{b}} \abs{f -
g}$（有限：$f - g$ 有界）；任一集合取\emph{离散}
距离（$d(x,y) = 1$，当 $x \neq y$ 时）。由范数导出的距离
将在\cref{ch:b2:nvs}中讨论。
\end{example}

\begin{definition}[度量空间的拓扑]\label{def:b2:metric:topology}
集合 $U \subseteq X$ 称为\emph{开集}\index{开集!度量空间中}，当
$U$ 中每一点都是某个含于 $U$ 的球的球心时；
$F$ 称为\emph{闭集}，当其补集为开集时。邻域、
内部、闭包、稠密性、边界的定义，与实直线上（本科一年级分册）
完全相同，只需以球代替区间；其中所证结论
——开集的并与交、内部与闭包的刻画、
闭包作为最小闭超集等结论——
在此照样成立，证明相同。开球是
开集，闭球是闭集（由三角不等式）。
\end{definition}

\begin{example}[同一集合上的内部、闭包与边界]\label{ex:b2:metric:intclosure}
在 $\R$ 中，设 $A = \intoc{0}{1} \cup \{2\}$。内部：
$\intoo{0}{1}$——对任一 $x \in \intoo01$，其附近的小球
都含于 $A$；对 $1$，每个球 $\intoo{1-r}{1+r}$ 都
在右侧溢出 $A$，故 $1$ 不是内点；而孤立点
$2$ 也不是内点。闭包：
$\intcc{0}{1} \cup \{2\}$（点 $0$ 是 $A$ 的极限，
其余不增添）。边界（闭包减去内部）：
$\{0, 1, 2\}$。值得记住的不对称性：一个
端点可以属于集合而非内点（$1$），也可以
是附着点却不属于集合（$0$）；孤立点
是它自身的边界（$2$）。同样的计算原样适用于
任一\omterm{def:b2:metric:def}{度量空间}，只需以球代替区间。
\end{example}

\begin{definition}[极限与连续]\label{def:b2:metric:continuity}
$x_n \to x$ 指在 $X$ 中 $d(x_n, x) \to 0$。\omterm{def:b2:metric:def}{度量空间}之间的映射 $f \colon X \to Y$
在 $a$ 处\emph{连续}\index{连续}，当
\[
\forall \varepsilon > 0,\ \exists\delta > 0,\quad
d_X(x, a) \leq \delta \implies d_Y\bigl(f(x), f(a)\bigr) \leq
\varepsilon ;
\]
等价地（证明与 $\R$ 上相同）：$f(x_n) \to f(a)$ 对每个
$x_n \to a$ 的序列都成立。$f$ 为具常数 $k$ 的\emph{利普希茨}\index{利普希茨}映射，当
$d_Y(f(x), f(y)) \leq k\, d_X(x, y)$ 恒成立，于是一致
连续，因而也连续。
\end{definition}

\begin{theorem}[连续性的整体刻画]\label{thm:b2:metric:globalcontinuity}
$f \colon X \to Y$ \omterm{def:b2:metric:continuity}{连续}（在每一点处）当且仅当
任一\omterm{def:b2:metric:topology}{开集}的原像是\omterm{def:b2:metric:topology}{开集}——等价地，
任一闭集的原像为闭集。
\end{theorem}

\begin{proof}
（$\Rightarrow$）设 $V \subseteq Y$ 为\omterm{def:b2:metric:topology}{开集}，$a \in f^{-1}(V)$：
则存在球 $B(f(a), \varepsilon) \subseteq V$；由 $a$ 处的\omterm{def:b2:metric:continuity}{连续}性，
存在 $\delta$ 使 $f(B(a, \delta)) \subseteq B(f(a),
\varepsilon)$，从而 $B(a, \delta) \subseteq f^{-1}(V)$。

（$\Leftarrow$）给定 $a$ 与 $\varepsilon$：$f^{-1}\bigl(B(f(a),
\varepsilon)\bigr)$ 是\omterm{def:b2:metric:topology}{开集}且含 $a$，故含一球
$B(a, \delta)$，这正是在 $a$ 处\omterm{def:b2:metric:continuity}{连续}的定义。闭集的情形
取补集即可（见\cref{prop:b2:structures:images}）。
\end{proof}

\section{完备空间}

\begin{definition}\label{def:b2:metric:complete}
称序列 $(x_n)$ 为\emph{柯西序列}，当 $\sup_{p, q \geq N} d(x_p,
x_q) \to 0$（当 $N \to \infty$ 时）。称\omterm{def:b2:metric:def}{度量空间}为
\emph{完备}\index{完备}，当每个柯西
序列都收敛。收敛 $\Rightarrow$ 柯西总成立；完备空间的闭子集
仍是完备的，任一空间中的完备子集都是闭集
（证明与 $\R$ 上相同，见本科一年级分册）。
\end{definition}

\begin{example}[无极限的柯西序列]\label{ex:b2:metric:cauchynolimit}
在 $X = \Q$ 上取通常距离，
$\sqrt2$ 的各级小数截断依次为
\[
x_0 = 1,\quad x_1 = 1.4,\quad x_2 = 1.41,\quad x_3 = 1.414,
\quad\dots
\]
满足 $\abs{x_p - x_q} \leq 10^{-\min(p,q)}$，故在 $\Q$ 中为柯西序列。
若 $\Q$ 中有极限，它也是 $\R$ 中的极限，即
$\sqrt2 \notin \Q$，故 $X$ 中不存在极限。不\omterm{def:b2:metric:complete}{完备}性即
存在这种“虚极限”。$\R$ 的\omterm{def:b2:metric:complete}{完备}性
正是在本科一年级分册中专门构造的，
使每个柯西序列都有归宿。
\end{example}

\begin{theorem}\label{thm:b2:metric:rncomplete}
$\R^n$（三种距离中的任一种，见
\cref{ex:b2:metric:examples}），以及 $\bigl(C(\intcc{a}{b}),
d_\infty\bigr)$ 都是\omterm{def:b2:metric:complete}{完备}的。
\end{theorem}

\begin{proof}
$\R^n$：柯西序列在每个坐标上仍为柯西序列（
$\abs{x_i - y_i} \leq d(x,y)$ 对三种距离均成立），于是每个
坐标都收敛（$\R$ 的\omterm{def:b2:metric:complete}{完备}性见本科一年级分册），并且
逐坐标收敛蕴含关于 $d_\infty$ 的收敛
（坐标个数有限），从而对三种距离都成立（三种距离彼此
在常数因子意义下互相控制：$d_\infty
\leq d_2 \leq d_1 \leq n\,d_\infty$）。

$C(\intcc{a}{b})$：设 $(f_n)$ 为关于 $d_\infty$ 的柯西序列。对每个 $x$，
$(f_n(x))$ 在 $\R$ 中为柯西序列（$\abs{f_p(x) - f_q(x)} \leq
d_\infty(f_p, f_q)$），故收敛于某个 $f(x)$。
在 $\abs{f_p(x) - f_q(x)} \leq \varepsilon$（对 $p, q
\geq N_\varepsilon$，所有 $x$ 成立）中令 $q \to \infty$ 取极限，得 $\abs{f_p(x) -
f(x)} \leq \varepsilon$ 对所有 $x$ 成立，亦即 $d_\infty(f_p, f) \leq
\varepsilon$，即一致收敛。该极限是\omterm{def:b2:metric:continuity}{连续}的：对给定的
$\varepsilon$，取 $p$ 使 $\sup\abs{f_p - f} \leq \varepsilon$，
再利用 $f_p$ 在 $a$ 处的\omterm{def:b2:metric:continuity}{连续}性以及下面的三项分解
\[
\abs{f(x) - f(a)} \leq \abs{f(x) - f_p(x)} + \abs{f_p(x) - f_p(a)}
+ \abs{f_p(a) - f(a)} \leq 3\varepsilon
\]
当 $x$ 充分接近 $a$ 时成立。（这种“$3\varepsilon$ 论证”
在\cref{ch:b2:funcseq}中以一致极限定理的形式再次出现。）
\end{proof}

\begin{example}[由连续性识别开集与闭集]\label{ex:b2:metric:recognize}
整体刻画（见
\cref{thm:b2:metric:globalcontinuity}）是日常处理拓扑问题时最常用的工具。
在 $\R^2$ 中，集合 $\{(x, y) : x^2 + y^2
< 1,\ y > x^3\}$ 为\omterm{def:b2:metric:topology}{开集}——它是 $g^{-1}(\intoo{-\infty}{1})
\cap h^{-1}(\intoo{0}{+\infty})$，其中\omterm{def:b2:metric:continuity}{连续}函数 $g(x,y) =
x^2 + y^2$ 与 $h(x, y) = y - x^3$，是两个开原像的交。
在 $\bigl(C(\intcc01), d_\infty\bigr)$ 中，满足
$f(0) = f(1)$ 且 $\int_0^1 f = 0$ 的函数全体所成集合为闭集——
它是 $\{(0,0)\}$ 在\omterm{def:b2:metric:continuity}{连续}映射 $f \mapsto
\bigl(f(0) - f(1),\ \int_0^1 f\bigr)$ 下的原像，映射取值于 $\R^2$，
其各分量均为 $1$-利普希茨的（见
\cref{exo:b2:metric:3}）。此法从不画图：先给出\omterm{def:b2:metric:continuity}{连续}映射，
将集合读作原像，再引用定理。
\end{example}

\begin{example}[由无穷多个条件定义的闭集]\label{ex:b2:metric:lipclosed}
在 $\bigl(C(\intcc01), d_\infty\bigr)$ 中，集合
\[
L = \{f : \abs{f(x) - f(y)} \leq \abs{x - y}
\ \text{对所有 } x, y\}
\]
为 $1$-利普希茨函数全体，是闭集，尽管它由
不可数多个条件界定：对每个固定数对 $(x, y)$，
映射 $f \mapsto \abs{f(x) - f(y)} - \abs{x - y}$ \omterm{def:b2:metric:continuity}{连续}
（求值映射是 $1$-利普希茨的），因此每个单独条件
都定义一个闭集，而 $L$ 是这族集合的\emph{交}：
任意闭集之交仍是闭集。同一模板
可用于证明单调函数、凸函数，
以及以固定 $g$ 为界的函数类的闭性。一致极限继承
一切可由一族逐点闭约束表达的性质，
而一致极限所
\emph{不}自动继承的性质——例如可微性——
这正是 \cref{ch:b2:funcseq} 必须下功夫处理的内容。
\end{example}

\begin{theorem}[巴拿赫不动点定理]\label{thm:b2:metric:banach}
设 $X$ 为非空\omterm{def:b2:metric:complete}{完备}\omterm{def:b2:metric:def}{度量空间}，$f \colon X \to X$ 为
\emph{压缩映射}，即常数 $k < 1$ 的\omterm{def:b2:metric:continuity}{利普希茨}映射。则 $f$ 有
唯一不动点 $\ell$，且每条轨道 $x_{n+1} = f(x_n)$
都收敛于 $\ell$，并有估计\index{巴拿赫不动点定理}：
\[
d(x_n, \ell) \leq \frac{k^n}{1 - k}\, d(x_1, x_0) .
\]
\end{theorem}

\begin{proof}
唯一性：两个不动点之间的距离 $\leq k$ 倍于其自身距离。
存在性：由归纳法得 $d(x_{n+1}, x_n) \leq k^n d(x_1, x_0)$，因此
对 $q > p$，
\[
d(x_q, x_p) \leq \sum_{j=p}^{q-1} d(x_{j+1}, x_j)
\leq d(x_1, x_0) \sum_{j \geq p} k^j
= \frac{k^p}{1-k}\, d(x_1, x_0) \xrightarrow[p\to\infty]{} 0 :
\]
柯西序列，故收敛于某个 $\ell$；$f$ 的\omterm{def:b2:metric:continuity}{连续}性使
$x_{n+1} = f(x_n)$ 在极限下成立，即 $\ell = f(\ell)$。误差界
即为令 $q \to \infty$ 时上式。
\end{proof}

\begin{example}[一个积分方程]\label{ex:b2:metric:integralequation}
取 $X = C(\intcc{0}{1})$（\omterm{def:b2:metric:complete}{完备}，
由\cref{thm:b2:metric:rncomplete}可知），考虑 $T(f)(x) = 1 + \frac12
\int_0^x f(t)\,\dd t$。对 $f, g \in X$：
\[
\abs{T(f)(x) - T(g)(x)} \leq \frac12 \int_0^x \abs{f - g} \leq
\frac12\, d_\infty(f, g),
\]
因此 $T$ 为 $\frac12$-压缩映射，它有唯一的\omterm{def:b2:metric:continuity}{连续}
不动点——即方程 $f' = \frac f2$、$f(0) = 1$ 的解，即
$\eu^{x/2}$。这一方案经推广后成为
在\cref{ch:b2:diffeq}中的柯西--利普希茨定理。
\end{example}

\begin{example}[数值不动点：$x = \cos x$]\label{ex:b2:metric:cosfix}
在\omterm{def:b2:metric:complete}{完备}\omterm{def:b2:metric:def}{度量空间} $X = \intcc{0}{1}$ 上，映射 $f = \cos$ 将 $X$
映入 $\intcc{\cos 1}{1} \subseteq X$，并且是压缩映射：由
中值不等式，
\[
\abs{\cos x - \cos y} \leq \bigl(\sup_{\intcc01}\abs{\sin}\bigr)
\abs{x - y} = (\sin 1)\abs{x - y},
\qquad \sin 1 \approx 0.841 < 1 .
\]
由巴拿赫不动点定理：$x = \cos x$ 在 $\intcc01$ 中有唯一解（从而在
$\R$ 中：任一实不动点位于 $\intcc{-1}{1}$，再作用一次后位于
$\intcc{\cos 1}{1}$ 中），且迭代
$x_{n+1} = \cos x_n$ 从任意初值出发都收敛到它：$x_\infty
\approx 0.739085$，即在计算器上反复按
余弦键所得的著名数字。误差界预言
$(\sin1)^n/(1 - \sin1)$ 的衰减，约每
$13$ 次按键多一位正确数字；本章周末问题
（第 14 题）的后验界可即时为每一步给出验证。
\end{example}

\section{紧性}

\begin{definition}\label{def:b2:metric:compact}
称\omterm{def:b2:metric:def}{度量空间} $X$ 为\emph{紧度量空间}\index{紧度量空间}，
若 $X$ 中每个数列都有子列\emph{在}
$X$ 中收敛（即布尔查诺–魏尔斯特拉斯性质）。子集称为紧的，
当它在诱导距离下为紧时。
\end{definition}

\begin{theorem}[基本性质]\label{thm:b2:metric:compactprops}
\begin{enumerate}
  \item 紧子集是闭集且有界；紧空间的闭子集
        是紧的。
  \item 在 $\R^n$ 中逆命题也成立：紧 $\iff$ 闭且
        有界。
  \item 紧空间的\omterm{def:b2:metric:continuity}{连续}像是紧的；
        非空紧空间上的\omterm{def:b2:metric:continuity}{连续}实函数
        有界，且可取到其上、下确界。
  \item （海涅）紧空间上的\omterm{def:b2:metric:continuity}{连续}映射一致
        \omterm{def:b2:metric:continuity}{连续}。
  \item 乘积：若 $X, Y$ 紧，则 $X \times Y$ 亦紧（度量取为
        $d\bigl((x,y),(x',y')\bigr) = d(x,x') + d(y,y')$）。
\end{enumerate}
\end{theorem}

\begin{proof}
（1）与本科一年级分册中的论证相同：一个趋于%
无穷远或收敛于外部的数列，都没有
在内部收敛的子列；第二个断言则在
所在的紧空间中抽取子列，并利用闭性。

（2）$\R^n$ 中的有界数列有逐分量收敛的
子列：先在第一个坐标上抽取子列，
（对 $\R$ 用布尔查诺–魏尔斯特拉斯定理），再从该子列出发在
第二个坐标上抽取，依此类推（共 $n$ 次\omterm{def:b2:metric:continuity}{连续}抽取）；闭性保证
极限仍落在该集合内。

（3）给定 $(f(x_n))$，抽取 $x_{\varphi(n)} \to x \in X$；
由\omterm{def:b2:metric:continuity}{连续}性得 $f(x_{\varphi(n)}) \to f(x) \in f(X)$。实值情形：
由 $f(X) \subseteq \R$ 的紧性，它是闭且有界的，并且
$\sup f(X) \in f(X)$（集合的上确界是其附着点，且
$f(X)$ 为闭集）。

（4）本科一年级分册的证明可逐字照搬，这里以度量
形式写出。设 $f \colon X \to Y$ 在紧空间 $X$ 上\omterm{def:b2:metric:continuity}{连续}，
但不一致\omterm{def:b2:metric:continuity}{连续}：存在 $\varepsilon > 0$，使得对每个
$n$，都能取到满足下式的点：
\[
d_X(x_n, y_n) \leq \frac{1}{n+1}
\qquad\text{且}\qquad
d_Y\bigl(f(x_n), f(y_n)\bigr) > \varepsilon .
\]
抽取 $x_{\varphi(n)} \to a \in X$；则 $y_{\varphi(n)} \to
a$ 亦然（两者距离趋于 $0$）。在 $a$ 处的\omterm{def:b2:metric:continuity}{连续}性
使两个像数列都趋于 $f(a)$，故
$d_Y\bigl(f(x_{\varphi(n)}), f(y_{\varphi(n)})\bigr) \to 0$——
与一致的间隔 $> \varepsilon$ 相矛盾。紧性
恰好提供了一件关键的东西：聚点 $a$，
在其处可施加普通的\omterm{def:b2:metric:continuity}{连续}性。

（5）先在 $X$ 的分量上抽取，然后再在
$Y$ 的分量上抽取。
\end{proof}

\begin{example}[海涅定理：有紧性与无紧性]\label{ex:b2:metric:heinework}
在 $\intcc{0}{1}$ 上，函数 $x \mapsto x^2$ 一致
\omterm{def:b2:metric:continuity}{连续}——海涅定理无需计算即可断言，但直接
估计更具启发性：
\[
\abs{x^2 - y^2} = \abs{x + y}\,\abs{x - y} \leq 2\abs{x - y},
\]
于是 $\delta = \varepsilon/2$ 对\emph{所有点
同时}成立。在 $\R$ 上，同一函数并非一致\omterm{def:b2:metric:continuity}{连续}：
取 $x_n = n$，$y_n = n + \frac1n$，间距 $\abs{x_n - y_n}
= \frac1n \to 0$，而 $\abs{x_n^2 - y_n^2} = 2 + \frac1{n^2}
\geq 2$：不存在单一的 $\delta$ 能对 $\varepsilon = 1$ 适用。
机理一目了然：局部\omterm{def:b2:metric:continuity}{利普希茨}常数 $\abs{x + y}$
在紧集上有界，而在 $\R$ 上无界——海涅
定理恰好断言：紧性能一致地控制这类
局部常数。
\end{example}

\begin{method}[证明一个集合是紧的]\label{met:b2:metric:compactness}
三种途径，按使用频率排列。（1）\emph{环境
判别：}在 $\R^n$ 中（或任一有限维
赋范空间，见\cref{ch:b2:nvs}），验证闭性——通常将其
视为原像，见\cref{ex:b2:metric:recognize}——并验证有界。（2）
\emph{继承性：}已知紧空间的闭子集是紧的；
有限个紧集的并与乘积是紧的；紧集的
\omterm{def:b2:metric:continuity}{连续}像是紧的。（3）\emph{直接
验证：}对任一数列抽取收敛子列——
通常借助逐坐标的\omterm{def:b2:metric:continuity}{连续}抽取。
要证明\emph{非紧性}，只需一个见证：
即一个没有收敛子列的数列，最常见的是
两两距离 $\geq \varepsilon$ 的点列。
\end{method}

\begin{example}[集合间的距离：紧性大显身手]\label{ex:b2:metric:setdistance}
在\omterm{def:b2:metric:def}{度量空间}中，设 $K$ 为紧集，$F$ 为闭集，且 $K \cap F = \emptyset$，
则
\[
d(K, F) = \inf\,\{d(x, y) : x \in K,\ y \in F\} > 0 :
\]
函数 $x \mapsto d(x, F)$ \omterm{def:b2:metric:continuity}{连续}
（见\cref{exo:b2:metric:11}），且在 $K$ 上为正（若 $d(x, F) = 0$，
则将使 $x \in \overline F = F$），因此它在
紧集 $K$ 上取得正的最小值
（由\cref{thm:b2:metric:compactprops}~(3)）。紧性并非
可有可无的装饰：对两个\emph{闭集}，下确界
可以为零而不被取到——在 $\R^2$ 中，双曲线 $F_1 = \{xy
= 1\}$ 与坐标轴 $F_2 = \{y = 0\}$ 是两个不相交的闭集，
且 $d(F_1, F_2) = 0$（点 $(n, \frac1n)$ 趋近于
坐标轴）。逃向无穷远，正是紧性所禁止的。
\end{example}

\begin{theorem}[博雷尔--勒贝格]\label{thm:b2:metric:borellebesgue}
\omterm{def:b2:metric:def}{度量空间} $X$ 是紧的，当且仅当 $X$ 的每个由
\omterm{def:b2:metric:topology}{开集}构成的覆盖都有\emph{有限}子覆盖。\index{博雷尔--勒贝格
定理}
\end{theorem}

\begin{proof}
（$\Leftarrow$）设 $(x_n)$ 没有收敛子列。下证
每个 $x \in X$ 都存在球 $B(x, r_x)$，使其中包含 $x_n$ 的
指标 $n$ 仅有有限多个。否则，每个球
$B(x, \frac1{k+1})$ 都含无穷多项，
于是可依次选取指标
\[
\varphi(0) < \varphi(1) < \varphi(2) < \cdots
\quad\text{使得}\quad
x_{\varphi(k)} \in B\Bigl(x, \frac{1}{k+1}\Bigr)
\]
（每步都能做到，正是因为仍有无穷多个
候选）便得一个收敛于 $x$ 的子列。
这些球 $B(x, r_x)$ 覆盖 $X$；若其中有限个覆盖了 $X$，
则指标集 $\N$ 将是有限个有限集的并，
矛盾。

（$\Rightarrow$）分两步。\emph{勒贝格数：}设
$(U_i)$ 为紧空间 $X$ 的开覆盖，则存在 $\rho > 0$，使每个
半径为 $\rho$ 的球都落在某个 $U_i$ 中。否则，对每个 $n$ 取
$x_n$，使 $B(x_n, \frac{1}{n+1})$ 不落在任何 $U_i$ 中；抽取
子列 $x_{\varphi(n)} \to x \in U_{i_0} \supseteq B(x, r)$；当 $n$ 充分大时，
$B(x_{\varphi(n)}, \frac{1}{\varphi(n)+1}) \subseteq B(x, r)
\subseteq U_{i_0}$：矛盾。\emph{全有界性：}对
每个 $\rho > 0$，有限个半径为 $\rho$ 的球覆盖 $X$。
否则归纳地取 $x_{n+1}$ 位于 $B(x_0, \rho) \cup
\dots \cup B(x_n, \rho)$ 之外，故该数列两两距离 $\geq
\rho$，故无柯西子列——从而无收敛子列：
矛盾。综上，以有限个球覆盖 $X$，球的半径为
$\rho$（勒贝格数），每个球都含于某个 $U_i$，即得有限
子覆盖。
\end{proof}

\begin{example}[$\varepsilon$-网的计数]\label{ex:b2:metric:epsnet}
全有界性（取自
\cref{thm:b2:metric:borellebesgue}的证明）在
$\intcc{0}{1}$ 上十分具体：对 $\varepsilon > 0$，$\lceil
\frac{1}{2\varepsilon}\rceil$ 个球，以 $\varepsilon,
3\varepsilon, 5\varepsilon, \dots$ 为中心、半径为 $\varepsilon$ 的
球即可覆盖它——约 $\frac1{2\varepsilon}$ 个，且任何覆盖都
不能用少于 $\frac{1}{2\varepsilon}$ 个球完成（每个
球覆盖的长度至多 $2\varepsilon$）。在 $\intcc01^2$ 中，
个数平方增长，量级为 $\varepsilon^{-2}$：覆盖数
按 $\varepsilon^{-d}$ 的量级增长，$d$ 为维数——这是紧性的一个定量
侧面，也解释了无穷维
的单位球（见\cref{ch:b2:nvs}，其中根本不存在有限
$\frac13$-网），因而不能是紧的。
\end{example}

\begin{example}[从覆盖读出紧性]\label{ex:b2:metric:covers}
半开区间 $\intoc{0}{1}$ 被\omterm{def:b2:metric:topology}{开集}
$U_n = \intoo{\frac1n}{2}$，$n \geq 1$ 所覆盖；任一有限子族都有
最大指标 $N$，且遗漏了 $\intoc{0}{\frac1N}$：没有有限
子覆盖，因此 $\intoc{0}{1}$ 不紧——这一点
序列式定义借 $x_n = \frac1n$ 看穿此点，其极限
$0$ 逃逸了。另一方面，添加单点 $0$
同时修复两种判断：在 $\intcc{0}{1}$ 上，任一此类覆盖
必含某个含 $0$ 的集合，它吞下了一整段初始
段，再用有限多个集合补齐其余。两种语言
（见\cref{thm:b2:metric:borellebesgue}）总是同成败——覆盖
恰好在序列察觉逃逸之处也察觉逃逸。
\end{example}

\begin{remark}[本分册内的几个视角]\label{rem:b2:metric:perspectives}
本章是本分册的承重墙；请留意每根支柱
在何处承重。\emph{\omterm{def:b2:metric:complete}{完备}性}：柯西
准则成为巴拿赫
空间中级数的收敛判别法（见\cref{ch:b2:series}），一致收敛在
\cref{ch:b2:funcseq}中，恰为\omterm{def:b2:metric:complete}{完备}
\omterm{def:b2:metric:def}{度量空间} $\bigl(C, d_\infty\bigr)$ 中的收敛，而柯西--利普希茨
（见\cref{ch:b2:diffeq}）即是以积分方程
形式出现的巴拿赫不动点定理。\emph{紧性}：它证明
范数的等价性（见\cref{ch:b2:nvs}），极值的
取得（用于\cref{ch:b2:diffcalc}中的优化），以及
最佳逼近的存在性（\cref{ch:b2:nvs}中的周末
问题）。\emph{连通性}：它使局部结论整体化
——微分方程解的唯一性，
曲线上的介值定理（见\cref{ch:b2:curves}），以及
$GL_n(\R)$ 的两个分支，定向理论
（见\cref{ch:b2:multint}）将它们分开。
\end{remark}

\begin{remark}[常见误区]\label{rem:b2:metric:pitfalls}
(i) “闭且有界蕴含紧”是关于
$\R^n$ 的定理，而非\omterm{def:b2:metric:def}{度量空间}的：一个无穷集
配以离散度量，在自身中是闭且有界的，但不紧
（见\cref{exo:b2:metric:4}），而
$C(\intcc01)$ 的闭单位球同样不成立（见\cref{ch:b2:nvs}）。(ii) \omterm{def:b2:metric:complete}{完备}性是
\emph{距离}的性质，而非拓扑的性质：$\R$
赋以 $d(x,y) = \abs{\arctan x - \arctan y}$ 后，仍有通常的
收敛数列，却不\omterm{def:b2:metric:complete}{完备}
（见\cref{exo:b2:metric:1}）。(iii) \omterm{def:b2:metric:continuity}{连续}双射未必
是同胚——圆周参数化（见
\cref{exo:b2:metric:7}）即为一例；定义域的紧性可以修正之。
(iv) 巴拿赫定理要求 $k < 1$ \emph{一致地}成立：
单凭条件 $d(f(x), f(y)) < d(x,y)$，在一个
非紧空间上什么也保证不了（见\cref{exo:b2:metric:5}）。(v) 连通
一般并不蕴含道路连通——但对于本书中遇到的
赋范空间的开子集，二者一致
（见\cref{ch:b2:nvs}）。
\end{remark}

\begin{remark}[本章的应用]
分析部分处处用到。$C(\intcc{a}{b})$ 的\omterm{def:b2:metric:complete}{完备}性
支撑了
在\cref{ch:b2:funcseq}中的收敛定理，以及
在\cref{ch:b2:diffeq}中的柯西–\omterm{def:b2:metric:continuity}{利普希茨}理论（本章
周末问题已证明了局部皮卡–林德勒夫定理）；紧性给出
有限维空间中范数的等价性（见\cref{ch:b2:nvs}），以及
在\cref{ch:b2:diffcalc}中极值的存在性；连通性为
在\cref{ch:b2:realfun}中的介值论证
及微分方程唯一性的整体化提供基础。在
本科三年级分册中，函数空间中的紧性（
阿尔泽拉–阿斯科利定理）与贝尔纲定理
（本章的\cref{exo:b2:metric:12}）成为日常使用的工具。
\end{remark}

\begin{omfigure}
\begin{tikzpicture}[xscale=6.4, yscale=0.8]
  \draw[omDef, line width=3pt] (0,3) -- (1,3);
  \draw[omDef, line width=3pt] (0,2) -- (0.3333,2);
  \draw[omDef, line width=3pt] (0.6667,2) -- (1,2);
  \foreach \a in {0, 0.6667} {
    \draw[omDef, line width=3pt] (\a,1) -- (\a + 0.1111,1);
    \draw[omDef, line width=3pt]
      (\a + 0.2222,1) -- (\a + 0.3333,1);
  }
  \foreach \a in {0, 0.2222, 0.6667, 0.8889} {
    \draw[omThm, line width=3pt] (\a,0) -- (\a + 0.037,0);
    \draw[omThm, line width=3pt]
      (\a + 0.0741,0) -- (\a + 0.1111,0);
  }
  \node[left, font=\small] at (0,3) {$C_0$};
  \node[left, font=\small] at (0,2) {$C_1$};
  \node[left, font=\small] at (0,1) {$C_2$};
  \node[left, font=\small] at (0,0) {$C_3$};
\end{tikzpicture}

{\small 康托尔集的前几个阶段
（见\cref{exo:b2:metric:8}）：每一级都删去每条线段中间的
开三分之一区间。$C = \bigcap_n C_n$ 是
紧集，内部为空、长度为零，却与
$\R$ \omterm{def:b2:structures:countable}{等势}——它又作为集合上某个压缩映射的\emph{不动点}
在本章周末问题（第
22 题）中再次出现。}
\end{omfigure}

\section{连通性}

\begin{definition}\label{def:b2:metric:connected}
$X$ 称为\emph{连通的}\index{连通空间}，当它不能
分解为两个非空开子集之并——等价地，
其仅有的既开又闭的子集是 $\emptyset$ 与 $X$。$X$ 称为
\emph{道路连通的}，当任意两点可由\omterm{def:b2:metric:continuity}{连续}
映射 $\gamma \colon \intcc{0}{1} \to X$ 相连。
\end{definition}

\begin{theorem}\label{thm:b2:metric:connectedness}
\begin{enumerate}
  \item $\R$ 的连通子集恰好是区间。
  \item \omterm{def:b2:metric:connected}{连通空间}的\omterm{def:b2:metric:continuity}{连续}像仍连通——
        由此得一般的介值定理：\omterm{def:b2:metric:connected}{连通空间}上的
        \omterm{def:b2:metric:continuity}{连续}实函数，在其任意两个函数值之间
        的每个值均可取到。
  \item 道路连通 $\Rightarrow$ 连通。（逆命题
        一般不成立，但对赋范空间中的开子集成立，
        见\cref{ch:b2:nvs}。）
\end{enumerate}
\end{theorem}

\begin{proof}
(1) 若 $A$ 不是区间，则它在其两点之间漏掉某个 $z$：
$A = (A \cap \intoo{-\infty}{z}) \cup (A \cap \intoo{z}{+\infty})$
上式将它分成两个非空的（相对于 $A$ 的）\omterm{def:b2:metric:topology}{开集}。反之，设
$I$ 为区间，$I = U \cup V$ 为其一个分解，两者均为非空的
相对\omterm{def:b2:metric:topology}{开集}；取 $a \in U$，$b \in V$，不妨设 $a < b$，并令
$s = \sup\,(U \cap \intcc{a}{b})$，它是 $\intcc{a}{b}
\subseteq I$ 中的一点。若 $s \in U$，则 $s \neq b$，且由
$U$ 的相对开性，$s$ 周围一个（与 $I$ 相交的）区间整个含于
$U$ 中，故 $U \cap \intcc{a}{b}$ 中有大于 $s$ 的点，
与上确界矛盾。若 $s \in V$：由 $V$ 的相对开性，
存在区间 $\intoo{s - r}{s + r} \cap I$ 含于 $V$；但上确界
是 $U \cap \intcc{a}{b}$ 的附着点，故该区间必与之相交，
与 $U \cap V = \emptyset$ 矛盾。（这正是一年级针对
$\R$ 的既开又闭论证，此处在 $I$ 内进行。）

(2) 若 $f(X) = U' \cup V'$ 分解为两个非空的相对\omterm{def:b2:metric:topology}{开
集}，则 $X = f^{-1}(U') \cup f^{-1}(V')$ 给出 $X$ 的一个分解
（见\cref{thm:b2:metric:globalcontinuity}）。介值定理：$f(X) \subseteq \R$
连通，故由 (1)，它是一个区间。

(3) 设 $X = U \cup V$，两者皆为非空\omterm{def:b2:metric:topology}{开集}，取 $a \in U$
与 $b \in V$，用道路 $\gamma$ 连接：则 $\gamma^{-1}(U),
\gamma^{-1}(V)$ 给出 $\intcc{0}{1}$ 的分解，与 (1) 矛盾。
\end{proof}

\begin{example}\label{ex:b2:metric:glnr}
$GL_n(\R)$ 不连通：$\det$ 是\omterm{def:b2:metric:continuity}{连续}的（矩阵元素的多项式），
满射到 $\R^*$ 上，后者不连通；$\R_+^*$ 与
$\R_-^*$ 的原像把 $GL_n(\R)$ 分成两块。（事实上每一块都
是道路\omterm{def:b2:metric:connected}{连通的}——这是超出本书需要的一个有趣练习。）与之相对，
$GL_n(\C)$ \emph{是}道路\omterm{def:b2:metric:connected}{连通的}：
见\cref{exo:b2:metric:10}。
\end{example}

\begin{example}[仅由连通性得到的不动点]\label{ex:b2:metric:ivtfixed}
每个\omterm{def:b2:metric:continuity}{连续}映射 $f \colon \intcc01 \to \intcc01$ 都有一个
不动点——无需压缩假设，也无需迭代。考虑
$g(x) = f(x) - x$，它在连通集 $\intcc01$ 上\omterm{def:b2:metric:continuity}{连续}：
\[
g(0) = f(0) \geq 0,
\qquad
g(1) = f(1) - 1 \leq 0 ,
\]
由介值定理
（见\cref{thm:b2:metric:connectedness}~(2)）给出
$g$ 的一个零点，即 $f$ 的一个不动点。与巴拿赫不动点定理
（见\cref{thm:b2:metric:banach}）相比：这里的存在性纯属拓扑性质，
且无代价，但唯一性与算法都已丧失——$f =
\mathrm{id}$ 的每一点都是不动点，而迭代一个
非压缩的 $f$ 可能永远循环。本章的两个不动点
定理回答的是不同的问题，
所用的“货币”也不相同。
\end{example}

\begin{example}[$\R$ 与 $\R^2$ 不同胚]\label{ex:b2:metric:rnotr2}
连通性是一种拓扑特征。假设 $h \colon
\R^2 \to \R$ 为同胚（即\omterm{def:b2:metric:continuity}{连续}双射，
逆映射\omterm{def:b2:metric:continuity}{连续}）。去掉一点 $a \in \R^2$，则
限制映射 $h \colon \R^2\setminus\{a\} \to
\R\setminus\{h(a)\}$ 仍为同胚。但 $\R^2$ 去掉
一点后是道路\omterm{def:b2:metric:connected}{连通的}——任意两点可用线段相连，
若直线段被 $a$ 挡住，则改走经过某辅助点的第二条线段，
因此连通
（\cref{thm:b2:metric:connectedness}~(3)）；而 $\R$ 去掉
一点后分成两个非空开半直线，故不连通。
连通性在\omterm{def:b2:metric:continuity}{连续}映射下保持不变，矛盾。
平面与直线作为\emph{
拓扑空间}确有本质区别——这一点仅凭基数
（\cref{exo:b2:structures:3}中那类双射确实存在！）太粗，
无法察觉。
\end{example}

\section{练习}

\begin{exercise}[$\star$]\label{exo:b2:metric:1}
在 $\R$ 上验证 $\delta(x, y) = \min(1, \abs{x - y})$ 与
$d(x,y) = \abs{\arctan x - \arctan y}$ 都是距离。
对每个距离，哪些数列收敛？$(\R, d)$ 是否\omterm{def:b2:metric:complete}{完备}？
\end{exercise}

\begin{exercise}[$\star$]\label{exo:b2:metric:2}
在\omterm{def:b2:metric:def}{度量空间}中，证明收敛数列是柯西列且
有界；并证明带收敛子列的柯西数列
收敛。再由此推出\omterm{def:b2:metric:compact}{紧度量空间}是\omterm{def:b2:metric:complete}{完备}的。
\end{exercise}

\begin{exercise}[$\star$]\label{exo:b2:metric:3}
在 $\bigl(C(\intcc{0}{1}), d_\infty\bigr)$ 中，计算
$f(x) = x$ 与 $g(x) = x^2$ 之间的距离；描述闭球
$\overline B(0, 1)$；并证明集合 $\{f : f(0) = 0\}$ 是
闭集，而 $\{f : f(0) > 0\}$ 是\omterm{def:b2:metric:topology}{开集}。
\end{exercise}

\begin{exercise}[$\star\star$]\label{exo:b2:metric:4}
证明离散\omterm{def:b2:metric:def}{度量空间} $X$（任意集合）是\omterm{def:b2:metric:complete}{完备}的，并
证明它紧当且仅当 $X$ 为有限集。哪些子集是
\omterm{def:b2:metric:connected}{连通的}？
\end{exercise}

\begin{exercise}[$\star\star$]\label{exo:b2:metric:5}
设 $X$ 为紧集，且 $f \colon X \to X$ 满足
\[
d\bigl(f(x), f(y)\bigr) < d(x, y) \quad \text{对所有 } x \neq y .
\]
证明 $f$ 有唯一不动点 \emph{（极小化 $x \mapsto
d(x, f(x))$）}，并在 $X = \intco{1}{+\infty}$
（非紧）上给出一个无不动点的例子。
\end{exercise}

\begin{exercise}[$\star\star$]\label{exo:b2:metric:6}
（嵌套紧集）设 $(K_n)$ 为一列递减的非空
\omterm{def:b2:metric:def}{度量空间}的紧子集。证明 $\bigcap_n K_n \neq
\emptyset$ \emph{（取 $x_n \in K_n$，再取子列）}。用例子说明，
$\R$ 中非空嵌套的\emph{闭集}
的交集可以为空。
\end{exercise}

\begin{exercise}[$\star\star$]\label{exo:b2:metric:7}
设 $K$ 为紧集，$f \colon K \to Y$ 为\omterm{def:b2:metric:continuity}{连续}双射。
证明 $f^{-1}$ \omterm{def:b2:metric:continuity}{连续} \emph{（利用闭集：
\cref{thm:b2:metric:globalcontinuity} 与
\cref{thm:b2:metric:compactprops}）}。去掉紧性，给出反例
（$\gamma(t) = (\cos t, \sin t)$ 在
$\intco{0}{2\pi}$ 上）。
\end{exercise}

\begin{exercise}[$\star\star$]\label{exo:b2:metric:8}
\emph{康托尔集} $C$ 由 $\intcc{0}{1}$ 出发，
反复删去中间的开三分之一区间。证明 $C$ 是紧的，
具有空内部，且是无限集——事实上与
$\{0,1\}^{\N}$ \omterm{def:b2:structures:countable}{等势} \emph{（三进制展开，数字取 $0,2$；
\cref{exo:b2:structures:3}）}。
\end{exercise}

\begin{exercise}[$\star\star\star$]\label{exo:b2:metric:9}
设 $X$ 为\emph{紧}\omterm{def:b2:metric:def}{度量空间}，$f \colon X \to X$ 为
等距映射：$d(f(x), f(y)) = d(x, y)$。证明 $f$ 是满射。
\emph{提示：若 $a \notin f(X)$，则 $\varepsilon = d(a, f(X)) > 0$
（为什么？）；考察轨道 $a, f(a), f^2(a), \dots$，并证明其
各点两两之间距离都 $\geq \varepsilon$，与
紧性矛盾。}
\end{exercise}

\begin{exercise}[$\star\star\star$]\label{exo:b2:metric:10}
证明 $GL_n(\C)$ 是道路\omterm{def:b2:metric:connected}{连通的}。\emph{提示：给定可逆矩阵 $A, B$，
考虑 $p(z) = \det\bigl((1 - z)A + zB\bigr)$，$z
\in \C$ 时：它是关于 $z$ 的多项式，不恒为零，故只有有限个根；
用一条道路连接 $0$ 与 $1$，
该道路位于 $\C$ 中且避开这些根。}
\end{exercise}

\begin{exercise}[$\star\star$]\label{exo:b2:metric:11}
对 $\emptyset \neq A \subseteq X$，令 $d(x, A) = \inf_{a \in A}
d(x, a)$。证明 $x \mapsto d(x, A)$ 为 $1$-利普希茨函数，
$d(x, A) = 0$ 当且仅当 $x \in \overline A$，并证明对不交的
非空\emph{闭集} $A, B$，函数
\[
\varphi(x) = \frac{d(x, A)}{d(x, A) + d(x, B)}
\]
是良定义的、\omterm{def:b2:metric:continuity}{连续}的，且取值为 $0$ 的点恰好是 $A$ 中的点，并
取值为 $1$ 的点恰好是 $B$ 中的点——一个分离任意两个
不交闭集的\omterm{def:b2:metric:continuity}{连续}“开关”。
\end{exercise}

\begin{exercise}[$\star\star\star$]\label{exo:b2:metric:12}
（贝尔）设 $X$ 为\omterm{def:b2:metric:complete}{完备}\omterm{def:b2:metric:def}{度量空间}，$(U_n)_{n\geq1}$ 为
一列稠密开子集。证明 $\bigcap_n U_n$ 在 $X$ 中
稠密\emph{（在任一球内，构造嵌套闭球
$\overline B(x_n, r_n) \subseteq U_n$，其中 $r_n \to 0$，并利用
\omterm{def:b2:metric:complete}{完备}性）}。由此推出 $\R$ 不是
内部为空的闭集的可数并，并再次推出
$\R$ 不可数。
\end{exercise}

\section{问题：皮卡迭代}

\omterm{def:b2:metric:complete}{完备}性加压缩是一台求解机器：把一个写成不动点问题的
方程交给它，它便返回
存在性、唯一性、算法以及误差界。这道周末问题
让这台机器全力运转，用于方程 $y' = f(t,
y)$：证明局部的\emph{皮卡--林德勒夫定理}（它是
在\cref{ch:b2:diffeq}中柯西--利普希茨理论的非线性核心），
逐一通过反例观察每个假设如何发挥作用，
并收获纯度量意义上的“红利”——
数据的\omterm{def:b2:metric:continuity}{连续}依赖性、开普勒方程，以及
康托集的自相似性。

\begin{omfigure}
\begin{tikzpicture}
\begin{axis}[omaxis, width=8.6cm, height=5.4cm,
    xmin=-0.05, xmax=1.45, ymin=0, ymax=4.2,
    xtick={0,1}, ytick={1,2,3}]
  \addplot[omDef!45, thick, domain=0:1.4, samples=2] {1};
  \addplot[omDef!65, thick, domain=0:1.4, samples=2] {1 + x};
  \addplot[omDef, thick, domain=0:1.4, samples=60]
    {1 + x + x^2/2};
  \addplot[omThm, very thick, domain=0:1.4, samples=120]
    {exp(x)};
  \node[font=\small, omDef] at (axis cs:1.28,0.8) {$y_0$};
  \node[font=\small, omDef] at (axis cs:1.28,2.15) {$y_1$};
  \node[font=\small, omDef] at (axis cs:1.28,2.95) {$y_2$};
  \node[font=\small, omThm] at (axis cs:1.17,3.85)
    {$\eu^{\,t}$};
\end{axis}
\end{tikzpicture}

{\small 对 $y' = y$，$y(0) = 1$ 的皮卡迭代：每经过
一次 $T(y)(t) = 1 + \int_0^t y$ 的作用便增加一个泰勒项，而
压缩性质使整个序列一致收敛到
$\eu^{\,t}$。}
\end{omfigure}

\begin{problem}[{周末问题——皮卡--林德勒夫
定理}]
\label{pb:b2:metric:1}
全文假设 $t_0 \in \R$，$y_0 \in \R$，$a, b > 0$，且 $f$ 是
矩形 $R = \intcc{t_0 - a}{t_0 +
a} \times \intcc{y_0 - b}{y_0 + b}$，其界为 $M =
\sup_R\,\abs f$，并且\emph{关于第二个变量为 $L$-利普希茨}：
$\abs{f(t, y) - f(t, z)} \leq L\abs{y - z}$，只要两点
都属于 $R$。令
\[
h = \min\Bigl(a, \frac bM\Bigr) \quad (\text{ 其中 } h = a
\text{ 若 } M = 0), \qquad I = \intcc{t_0 - h}{t_0 + h}.
\]

\medskip
\noindent\textbf{第一部分——\omterm{def:b2:metric:complete}{完备}阶段。}
\begin{enumerate}
  \item 证明
        在\cref{def:b2:metric:complete}中引用的两个结论：\omterm{def:b2:metric:complete}{完备}\omterm{def:b2:metric:def}{度量空间}的闭子集
        是\omterm{def:b2:metric:complete}{完备}的，而\omterm{def:b2:metric:complete}{完备}子集
        在任意\omterm{def:b2:metric:def}{度量空间}中都是闭的。由此推出
        $\bigl(C(I), d_\infty\bigr)$ 的每个闭子集
        都是\omterm{def:b2:metric:complete}{完备}\omterm{def:b2:metric:def}{度量空间}。
  \item 证明映射 $C(I) \to C(I)$，$y \mapsto \bigl(t
        \mapsto \int_{t_0}^{t}y(s)\,\dd s\bigr)$，是
        $h$-利普希茨的（关于 $d_\infty$）。
  \item （不动点的移动量小于映射的移动量）设 $g \colon X
        \to X$ 为\omterm{def:b2:metric:def}{度量空间}上的 $k$-压缩映射，
        其不动点为 $\ell_g$；设 $\widetilde g \colon X \to X$
        为\emph{任一}具有不动点 $\ell_{\widetilde g}$ 的映射。
        证明
        \[
        d(\ell_g, \ell_{\widetilde g}) \leq
        \frac{d\bigl(g(\ell_{\widetilde g}),
        \widetilde g(\ell_{\widetilde g})\bigr)}{1 - k}
        \leq \frac{\sup_{x \in X} d\bigl(g(x), \widetilde
        g(x)\bigr)}{1 - k}.
        \]
  \item （迭代技巧）设 $X$ 为非空\omterm{def:b2:metric:complete}{完备}\omterm{def:b2:metric:def}{度量空间}，$g
        \colon X \to X$ 为映射——不假设它\omterm{def:b2:metric:continuity}{连续}——
        使得某个迭代 $g^m$ 为 $k$-压缩映射。
        证明 $g$ 有唯一不动点 $\ell$，且
        \emph{每一条}轨道 $x_{n+1} = g(x_n)$ 收敛到
        $\ell$。\emph{（$g$ 的不动点都是
        $g^m$ 的不动点；反之 $g(\ell)$ 是 $g^m$ 的不动点；
        沿模 $m$ 的剩余类分割轨道。）}
\end{enumerate}

\medskip
\noindent\textbf{第二部分——皮卡--林德勒夫定理。}
\begin{enumerate}[resume]
  \item 证明：函数 $y \colon I \to \intcc{y_0 - b}{y_0
        + b}$ 是 $C^1$ 的，且 $y(t_0) = y_0$，$y' = f(t, y)$ 在
        $I$ 上，当且仅当它\omterm{def:b2:metric:continuity}{连续}并满足
        积分方程
        \[
        y(t) = y_0 + \int_{t_0}^{t} f\bigl(s, y(s)\bigr)\dd s
        \qquad (t \in I).
        \]
  \item 设 $X_h = \{y \in C(I) : \abs{y(t) - y_0} \leq b
        \text{ 在 } I\}$，并令 $T$ 由 $T(y)(t) =
        y_0 + \int_{t_0}^{t}f(s, y(s))\dd s$ 定义。证明 $X_h$
        是 $C(I)$ 的非空闭子集，从而\omterm{def:b2:metric:complete}{完备}，并且
        $T$ 把 $X_h$ 映入 $X_h$——这正是 $h \leq
        b/M$ 起作用的地方。
  \item 证明 $d_\infty\bigl(T(y), T(z)\bigr) \leq
        Lh\,d_\infty(y, z)$（在 $X_h$ 上）：若 $Lh < 1$，
        则巴拿赫不动点定理已足以给出结论。我们将
        在下面去掉这一“足够小”的条件。
  \item 对 $n$ 用归纳法证明：
        \[
        \abs{T^n(y)(t) - T^n(z)(t)} \leq
        \frac{\bigl(L\abs{t - t_0}\bigr)^n}{n!}\,
        d_\infty(y, z) \qquad (y, z \in X_h,\ t \in I),
        \]
        于是 $T$ 的某个迭代为压缩映射。利用
        第 4 小题（\emph{皮卡--林德勒夫定理}）得出：柯西问题
        $y' = f(t,y)$，$y(t_0) = y_0$ 恰有
        一个解，它定义在 $I = \intcc{t_0 - h}{t_0 + h}$ 上，
        且取值于 $\intcc{y_0 - b}{y_0 + b}$ 中。
  \item 证明限制条件“取值于 $\intcc{y_0 -
        b}{y_0 + b}$”是自动成立的：任一
        柯西问题的解，若其定义在 $I$ 上且图像起于 $R$，则
        其取值始终位于 $\intcc{y_0 - b}{y_0 + b}$ 中\emph{（考虑首次离开时刻，
        并估计 $\abs{y(t) - y_0}$，以 $M\abs{t -
        t_0}$ 为界）}。由此，唯一性在所有定义于
        $I$ 上的解中成立。
  \item 对 $y' = y$，$y(0) = 1$ 运行这台机器，从
        常值函数 $y^{(0)} \equiv 1$ 出发：计算皮卡迭代
        $y^{(n)}$，辨认出它们，并描述
        收敛情况。
\end{enumerate}

\medskip
\noindent\textbf{第三部分——每个假设都有其用处。}
\begin{enumerate}[resume]
  \item （\omterm{def:b2:metric:continuity}{利普希茨}条件不成立，唯一性不成立）对 $y' =
        2\sqrt{\abs y}$，$y(0) = 0$，验证 $y \equiv 0$
        以及对每个 $c \geq 0$，函数 $y_c(t) = 0$ 在
        $t \leq c$ 上，$y_c(t) = (t - c)^2$ 在 $t > c$ 上，皆为
        $C^1$ 解，定义于 $\R$ 上。映射 $y \mapsto
        2\sqrt{\abs y}$ 究竟在何处不满足\omterm{def:b2:metric:continuity}{利普希茨}条件？
  \item （局部性是真实的）对 $y' = y^2$，$y(0) = 1$，
        显式求解，给出最大存在区间，并
        计算定理能够保证的最优 $h$，
        在所有矩形选取（$a$ 充分大，$b$ 任意）下，证明
        $h_{\max} = \sup_{b>0} \frac{b}{(1+b)^2} = \frac14$,
        而真解定义在 $\intoo{-\infty}{1}$ 上。
  \item （\omterm{def:b2:metric:complete}{完备}性并非摆设）在 $X = \Q \cap
        \intcc{1}{2}$ 上取通常的距离，令 $g(x) =
        \frac x2 + \frac1x$。证明 $g(X) \subseteq X$，且 $g$
        是 $\frac12$-压缩映射 \emph{（中值
        不等式）}，并证明 $g$ 在 $X$ 中没有不动点。
        指出巴拿赫定理的哪一条假设不成立，以及
        \omterm{def:b2:metric:complete}{完备}化空间中的不动点是什么？
  \item （误差估计）对\omterm{def:b2:metric:complete}{完备}空间上的 $k$-压缩映射 $g$，
        证明\emph{后验估计} $d(x_n,
        \ell) \leq \frac{k}{1-k}\,d(x_n, x_{n-1})$。对海伦
        映射 $g(x) = \frac x2 + \frac 1x$，定义于 $\intcc{1}{2}$ 上，
        （不动点 $\sqrt2$），取初值 $x_0 = \frac32$：
        按\emph{先验界}，
        $\frac{k^n}{1-k}d(x_1, x_0)$ 要使精度达到
        $10^{-6}$，需要多少步？实际上多少步已足够？
        （计算 $x_1, x_2, x_3$ 及其误差；压缩界
        可靠但偏于悲观——
        海伦迭代是二次收敛的。）
\end{enumerate}

\medskip
\noindent\textbf{第四部分——\omterm{def:b2:metric:continuity}{连续}依赖性。}本部分中
        $Lh < 1$，故 $T$ 本身是 $X_h$ 上的压缩映射
        （第 7 题）。记 $y[\,y_0\,]$ 表示初值为
        $y_0$ 的解。
\begin{enumerate}[resume]
  \item （对初值的依赖性）设 $z_0$ 为另一初值，
        满足 $\abs{z_0 - y_0}$ 充分小，使两个问题
        都落在矩形内。利用第 3 题，
        证明
        \[
        d_\infty\bigl(y[y_0], y[z_0]\bigr) \leq
        \frac{\abs{y_0 - z_0}}{1 - Lh}.
        \]
  \item （逐段接续得到长区间）假设解存在于一条长线段上，
        该线段被切分为 $m$ 段相邻的小段，
        每一段上前述估计在 $Lh \leq
        \frac12$ 时成立。证明每一段上偏差至多放大
        $2$ 倍，从而整体有 $d_\infty \leq 2^m\abs{y_0 - z_0}$
        整体而言——一个随长度指数增长的界，正是
        格朗沃尔引理中 $\eu^{L\abs{t - t_0}}$ 的离散对应
        （见\cref{ch:b2:diffeq}）。
  \item （对向量场的依赖性）设 $g$ 为另一向量场，
        在 $R$ 上同样 $L$-利普希茨（关于 $y$），且 $\sup_R \abs{f - g}
        \leq \varepsilon$。证明对应的
        解满足 $d_\infty \leq
        \frac{\varepsilon h}{1 - Lh}$：建模误差
        呈线性传播。
  \item （参数）若向量场族 $f_\lambda$
        一致地 $L$-利普希茨（关于 $y$），且 $\sup_R\abs{f_\lambda
        - f_\mu} \leq C\abs{\lambda - \mu}$，由此推出
        $\lambda \mapsto y_\lambda$ 是从
        参数空间到 $\bigl(C(I), d_\infty\bigr)$ 的\omterm{def:b2:metric:continuity}{利普希茨}映射。
  \item （方程组并不增加难度）说明为何第一、二、四部分
        对 $y$ 取值于 $\R^n$ 的情形逐字成立（上确界
        距离可由
        \cref{ex:b2:metric:examples}中的任一距离构造），然后计算全部皮卡
        迭代，对系统 $y' = Ay$，$y(0) = (c_1,
        c_2)$, $A = \begin{psmallmatrix} 0 & 1\\ 0 &
        0\end{psmallmatrix}$：证明迭代在一步后
        即稳定于精确解。
\end{enumerate}

\medskip
\noindent\textbf{第五部分——度量的收获与综合。}
\begin{enumerate}[resume]
  \item （恒等映射的扰动）设 $X$ 为\omterm{def:b2:metric:complete}{完备}的
        类赋范空间，取 $X = C(I)$ 或 $\R^n$。若
        $\eta \colon X \to X$ 是 $k$-利普希茨的，且 $k < 1$，
        证明 $x \mapsto x + \eta(x)$ 是 $X$ 上的双射，
        其逆映射是 $\frac1{1-k}$-利普希茨的 \emph{（对每个
        $y$，对 $x \mapsto y - \eta(x)$ 应用巴拿赫定理）}。这是
        反函数定理的度量核心
        （见\cref{ch:b2:diffcalc}）。
  \item （开普勒方程）对 $0 \leq e < 1$ 与 $m \in \R$，
        证明 $x = m + e\sin x$ 恰有一解，
        且迭代 $x_{n+1} = m + e\sin x_n$ 从任意起点出发
        都收敛于它，并作估计：对 $e = \frac12$，
        $m = 1$，由先验界，需多少次迭代方能保证误差 $\leq
        10^{-3}$？（解为 $x
        \approx 1.4987$。）
  \item （康托尔集是不动点）设 $S_1(x) = \frac
        x3$，$S_2(x) = \frac x3 + \frac23$，二者均定义于 $\R$ 上，并设
        $C$ 为\cref{exo:b2:metric:8}中的康托尔集。证明
        $C = S_1(C) \cup S_2(C)$，并用一句话说明
        为何没有\emph{其他}非空紧集满足此方程
        （映射 $A \mapsto S_1(A) \cup S_2(A)$ 在紧集之间某一距离下是
        压缩映射——该距离即豪斯多夫距离，
        本科三年级分册中将严格讨论）。
  \item （连通性使唯一性成为整体性质）设 $f$ 为\omterm{def:b2:metric:topology}{开集}上
        关于 $y$ 局部\omterm{def:b2:metric:continuity}{利普希茨}的，设 $y, z$ 为
        $y' = f(t, y)$ 在公共区间 $J$ 上的两个解，
        满足 $y(t_0) = z(t_0)$。证明 $y = z$ 在 $J$ 上，只需证
        $\{t \in J : y(t) = z(t)\}$ 非空、在 $J$ 中闭，
        且在 $J$ 中开（由局部唯一性），并利用区间的
        连通性
        （\cref{thm:b2:metric:connectedness}）。
  \item （线性方程无需小量假设）对 $y' =
        \alpha(t)y + \beta(t)$，其中 $\alpha, \beta$ 在
        线段 $\intcc{A}{B}$ 上\omterm{def:b2:metric:continuity}{连续}，改编第 8 题以证明
        阶乘型估计在\emph{整个}线段上成立，
        从而存在性与唯一性在该处整体成立——这是
        \cref{ch:b2:diffeq}中柯西–\omterm{def:b2:metric:continuity}{利普希茨}定理的标量情形，
        不对区间长度 $B - A$ 作任何限制。
  \item （综合）每项各用一句话说明：\omterm{def:b2:metric:complete}{完备}性贡献了什么；
        压缩性贡献了什么；
        与朴素的巴拿赫定理相比，迭代技巧带来了什么；
        连通性在何处起作用；以及第三部分中的哪个反例
        保护哪一条假设。写出最值定理的名称，
        并说明当 $f$ 仅为\omterm{def:b2:metric:continuity}{连续}时，用什么代替压缩性
        （皮亚诺定理，借助函数空间中的紧性——本科三年级分册中的
        阿尔泽拉–阿斯科利定理）。
\end{enumerate}
\end{problem}
